\documentclass{article}

% Language setting
\usepackage[english]{babel}

% Set page size and margins
\usepackage[letterpaper,top=2cm,bottom=2cm,left=3cm,right=3cm,marginparwidth=1.75cm]{geometry}

% Useful packages
\usepackage{amsmath}
\usepackage{amssymb}
\usepackage{graphicx}
\usepackage[colorlinks=true, allcolors=black]{hyperref}
\usepackage{titlesec}
\usepackage{xparse}

\NewDocumentCommand{\qcq}{m o o m}{%
	\ensuremath{\forall #1 \IfValueT{#2}{, #2} \IfValueT{#3}{, #3}  \in  \mathbb{#4}}%
}
\NewDocumentCommand{\ext}{m o o m}{%
	\ensuremath{\exists #1 \IfValueT{#2}{, #2} \IfValueT{#3}{, #3}  \in  \mathbb{#4}}%
}

\title{\textbf{Suites Numériques}}
\author{\textbf{Yassin Akkab}}

\newcommand{\R}{\mathbb{R}}
\newcommand{\N}{\mathbb{N}}
\newcommand{\Z}{\mathbb{Z}}
\newcommand{\C}{\mathbb{C}}
\newcommand{\Q}{\mathbb{Q}}
\newcommand{\D}{\mathbb{D}}

\begin{document}
	
	\maketitle
	\tableofcontents
	\newpage
	
	\section{Généralités sur les suites numériques}
	
	\subsection{Définition}
	Une suite numérique $u$ est une fonction définie sur une partie $I$ de $\mathbb{N}$ (généralement $\mathbb{N}$ ou $\mathbb{N} \setminus \{0, \dots, n_0\}$) à valeurs dans $\mathbb{R}$.
	On note $u_n$ l'image d'un entier $n$ par $u$, appelée le terme général de la suite. La suite elle-même est notée $(u_n)_{n \in I}$.
	
	\subsection{Majorabilité, minorabilité et suites bornées}
	Soit $(u_n)_{n \ge n_0}$ une suite numérique.
	\begin{itemize}
		\item La suite $(u_n)$ est dite \textbf{majorée} s'il existe un réel $M$ tel que :
		\[ \qcq{n}[\ge n_0]{N}, \quad u_n \le M \]
		\item La suite $(u_n)$ est dite \textbf{minorée} s'il existe un réel $m$ tel que :
		\[ \qcq{n}[\ge n_0]{N}, \quad u_n \ge m \]
		\item La suite $(u_n)$ est dite \textbf{bornée} si elle est à la fois majorée et minorée. Cela équivaut à :
		\[ \ext{A}[\ge 0]{R}, \quad \qcq{n}[\ge n_0]{N}, \quad |u_n| \le A \]
	\end{itemize}
	
	\subsection{Monotonie d'une suite}
	Une suite $(u_n)_{n \ge n_0}$ est dite :
	\begin{itemize}
		\item \textbf{Croissante} si : $\qcq{n}[\ge n_0]{N}, \quad u_{n+1} \ge u_n$
		\item \textbf{Décroissante} si : $\qcq{n}[\ge n_0]{N}, \quad u_{n+1} \le u_n$
		\item \textbf{Monotone} si elle est croissante ou décroissante.
	\end{itemize}
	
	\section{Limites et convergence d'une suite}
	
	\subsection{Suite convergente}
	Une suite $(u_n)_{n \ge n_0}$ converge vers un réel $l$ si tout intervalle ouvert contenant $l$ contient tous les termes de la suite à partir d'un certain rang. Formellement :
	\[ \lim_{n \to +\infty} u_n = l \iff \qcq{\varepsilon}[> 0]{R}, \ \ext{N}[\ge n_0]{N}, \ \qcq{n}[\ge N]{N}, \ |u_n - l| < \varepsilon \]
	Une suite qui n'est pas convergente est dite \textbf{divergente} (elle tend vers $\pm \infty$ ou n'admet pas de limite).
	
	\subsection{Propriétés des limites}
	\begin{itemize}
		\item \textbf{Unicité de la limite} : Si une suite converge, sa limite est unique.
		\item Toute suite convergente est bornée. (La réciproque est fausse, ex: $u_n = (-1)^n$).
		\item Si $(u_n)$ et $(v_n)$ convergent vers $l$ et $l'$ respectivement, alors :
		\begin{itemize}
			\item $\lim (u_n + v_n) = l + l'$
			\item $\lim (u_n \cdot v_n) = l \cdot l'$
			\item Si $l' \ne 0$, $\lim \frac{u_n}{v_n} = \frac{l}{l'}$
		\end{itemize}
	\end{itemize}
	
	\subsection{Théorèmes de convergence}
	\begin{itemize}
		\item \textbf{Théorème des limites monotones} :
		\begin{itemize}
			\item Toute suite croissante et majorée converge.
			\item Toute suite décroissante et minorée converge.
		\end{itemize}
		\item \textbf{Théorème des gendarmes} :
		Si $\qcq{n}[\ge n_0]{N}$, $v_n \le u_n \le w_n$ et $\lim v_n = \lim w_n = l$, alors la suite $(u_n)$ converge et $\lim u_n = l$.
	\end{itemize}
	
	\section{Suites adjacentes}
	Deux suites $(u_n)$ et $(v_n)$ sont dites \textbf{adjacentes} si :
	\begin{enumerate}
		\item L'une est croissante et l'autre est décroissante.
		\item $\lim_{n \to +\infty} (v_n - u_n) = 0$.
	\end{enumerate}
	\textbf{Théorème} : Si deux suites $(u_n)$ et $(v_n)$ sont adjacentes (avec $(u_n)$ croissante et $(v_n)$ décroissante), alors elles convergent vers la même limite $l$, et on a :
	\[ \qcq{n}{N}, \quad u_n \le l \le v_n \]
	
	\section{Suites récurrentes du type $u_{n+1} = f(u_n)$}
	Soit $f$ une fonction numérique continue sur un intervalle $I$ de $\mathbb{R}$, et $(u_n)$ la suite définie par $u_0 \in I$ et :
	\[ \qcq{n}{N}, \quad u_{n+1} = f(u_n) \]
	\textbf{Théorème (du point fixe)} :
	Si la suite $(u_n)$ vérifie les conditions suivantes :
	\begin{enumerate}
		\item $f(I) \subset I$ (l'intervalle $I$ est stable par $f$),
		\item La fonction $f$ est continue sur $I$,
		\item La suite $(u_n)$ converge vers une limite $l \in I$,
	\end{enumerate}
	alors $l$ est un point fixe de $f$, c'est-à-dire :
	\[ f(l) = l \]
	
\end{document}
